\scanpage{51}

\mspage{- 134 -}\hfill n$^\circ$ 371

\S7. \uuline{Algèbre extérieure d'un espace vectoriel.}

1. \uline{Application des déterminants aux équations linéaires sur un corps commutatif.}

Les résultats des \S\S\ 5 et 6 donnent en particulier les suivants pour les
espaces vectoriels :

PROPOSITION 1.- \uline{Soit} $E$ \uline{un espace vectoriel sur un corps commutatif} ; \uline{pour que} $p$
\uline{vecteurs} $x_i$ ($1\leq i\leq p$) \uline{de} $E$ \uline{forment un système libre, il faut et il suffit que}
\[
x_1\wedge x_2\wedge\ldots\wedge x_p\ne0.
\]

C'est un cas particulier du \S6, n$^\circ$5, prop.7, tout élément $\ne0$ d'un corps étant
inversible.

COROLLAIRE 1.- \uline{Soit} $\underline X$ \uline{une matrice de type} $(m,n)$ \uline{sur un corps commutatif. Le rang
de} $\underline X$ \uline{est égal au plus grand des entiers} $p$ \uline{tels qu'il existe au moins un mineur
d'ordre} $p$ \uline{de} $\underline X$ \uline{qui ne soit pas nul.}

En effet, $\operatorname{rg}(\underline X)$ est le nombre maximum des colonnes de $\underline X$ linéairement indépendantes
(chap.II, 3e éd., \S10, n$^\circ$\quad ) ; le corollaire résulte donc de la prop.1 et de
l'expression des composantes d'un $p$-vecteur décomposable (\S6, n$^\circ$3, prop.5).

COROLLAIRE 2.- \uline{Pour qu'une matrice carrée sur un corps commutatif soit inversible,
il faut et il suffit que son déterminant ne soit pas nul.}

COROLLAIRE 3.- \uline{Pour qu'un système linéaire} homogène \uline{de} $n$ \uline{équations à} $n$ \uline{inconnues
sur un corps commutatif admette une solution non triviale, il faut et il suffit que
le déterminant de sa matrice soit nul.}

Considérons un système de $m$ équations linéaires à $n$ inconnues sur un corps commutatif $K$ :
\[
(1)\qquad\sum_{j=1}^n\alpha_{ij}\xi_j=\beta_i\qquad(1\leq i\leq m).
\]
La matrice $\underline A=(\alpha_{ij})$ de ce système est donc de type $(m,n)$. Soit $\underline B$ la matrice
de type $(m,n+1)$ obtenue en \uline{bordant} $\underline A$ par la $(n+1)$-ème colonne $(\beta_i)$ ; on sait

\margin{\rotatebox{180}{50}}
